% Generated draft: requires mathematical and source review.
\documentclass[11pt]{article}
\usepackage[margin=25mm]{geometry}
\usepackage{fontspec}
\setmainfont{Latin Modern Roman}
\usepackage{amsmath,amssymb,mathtools,unicode-math}
\setmathfont{Latin Modern Math}
\usepackage{enumitem,needspace,xurl,hyperref}
\setlength{\parindent}{0pt}
\setlength{\parskip}{5pt}
\setlength{\emergencystretch}{3em}
\newcommand{\N}{\mathbb{N}}
\newcommand{\Z}{\mathbb{Z}}
\newcommand{\Q}{\mathbb{Q}}
\newcommand{\R}{\mathbb{R}}
\newcommand{\C}{\mathbb{C}}
\newcommand{\F}{\mathbb{F}}
\newcommand{\A}{\mathbb{A}}
\newcommand{\PP}{\mathbb{P}}
\newcommand{\E}{\mathbb{E}}
\newcommand{\eps}{\varepsilon}
\newcommand{\dd}{\,\mathrm{d}}
\begin{document}
\section*{UnsolvedMath ID: 701053}
\textbf{Source problem code:} catalogue entry
\subsection{Definitions and mathematical statement}
For each integer \(n>2\), let \(R_n=\C[x_1,\ldots,x_n]\). Write \(\operatorname{Aut}_{\C}(R_n)\) for the group of \(\C\)-algebra automorphisms, which fix every constant in \(\C\). A linear automorphism means an automorphism \(L\) for which there is a matrix \(A=(a_{ij})\in\operatorname{GL}_n(\C)\) satisfying
\[
L(x_i)=\sum_{j=1}^{n}a_{ij}x_j\qquad(1\leq i\leq n).
\]
An automorphism \(F\) has order exactly \(2\) if \(F^2=\operatorname{id}_{R_n}\) and \(F\neq\operatorname{id}_{R_n}\).

Under the interpretation that the source means \(\C\)-algebra automorphisms, the question is: for every integer \(n>2\) and every \(F\in\operatorname{Aut}_{\C}(R_n)\) of order exactly \(2\), do there exist \(G\in\operatorname{Aut}_{\C}(R_n)\) and \(A=(a_{ij})\in\operatorname{GL}_n(\C)\) such that
\[
(G^{-1}\circ F\circ G)(x_i)=\sum_{j=1}^{n}a_{ij}x_j\qquad(1\leq i\leq n)?
\]
The source says only ``automorphism of the polynomial ring'' and does not explicitly require constants to be fixed. Consequently, the choice between ring automorphisms and \(\C\)-algebra automorphisms remains a qualification requiring review; the displayed formulation is conditional on the latter interpretation. Here ``linear'' is interpreted as homogeneous linear substitution.

The supplied source calls this a well-known problem and states that the assertion is true for \(n=2\). That boundary-case assertion is retained as supplied background, without independent verification. The registry labels the problem open; this is a supplied status label, not an independently verified mathematical assessment.
\subsection{Short English statement}
Under the \(\C\)-algebra interpretation, must every automorphism of exact order \(2\) of \(\C[x_1,\ldots,x_n]\), for every \(n>2\), be conjugate to a linear automorphism? The source states that the answer is yes for \(n=2\), but does not explicitly specify whether its automorphisms fix \(\C\).
\subsection{Sources}
\begin{itemize}
\item Local UnsolvedMath registry, permanent ID \(701053\), exact source code \texttt{catalogue entry}, title ``Kourovka Notebook Problem 14.68.'' The supplied record labels the problem open.
\item \emph{The Kourovka Notebook (Unsolved Problems in Group Theory)}, original problem 14.68, 14th issue (1999); attributed to M. V. Neshchadim and described in the supplied statement as a well-known problem.
\item The supplied background identifies the 21st issue edition (2026), \texttt{arXiv:1401.0300v48}. The supplied collection description names E. I. Khukhro and V. D. Mazurov as editors. These bibliographic details and the assertion for \(n=2\) are supplied source information, not independently verified here.
\end{itemize}
\end{document}
